% 第5章 アインシュタイン方程式の線形化（WP5 skeleton; 本文・導出は執筆予定）
\section{アインシュタイン方程式の線形化}
\label{ch:ch05_einstein}

\paragraph{この章で導くこと（入力→出力）}
Newtonianゲージで $\delta G^\mu{}_\nu=8\pi G\,\delta T^\mu{}_\nu$ を成分計算し、実装に使う2本（$\Phi$ 発展式・$\Psi$ 代数式）に整理する。

% --- 節構成（01_textbook_spec.md より）-------------------------------------
\subsection{摂動 Christoffel・Ricci}
% TODO(WP5): 5.1 の本文。
\subsection{$\delta T^\mu{}_\nu$ と異方性応力（契約 S5.2）}
% TODO(WP5): 5.2 の本文。
\subsection{4本の方程式と独立2本の選択}
% TODO(WP5): 5.3 の本文。
\subsection{ニュートン極限の確認}
% TODO(WP5): 5.4 の本文。


\paragraph{導出契約}
本章の導出契約: S5.1–S5.2。各契約は「入力式→出力式」を中間ステップ省略なし
（1ステップ＝学部4年生が5分で検算可能）で履行する。\emph{本文プローズは WP5 で執筆。}

\paragraph{まとめ・チェックリスト・演習}
% TODO(WP5): 章末まとめ、チェックリスト、演習 3–6 問（うち1問は対応 NB の課題）。
